\documentclass[10pt]{article}

% graphicx package, useful for including eps, pdf, and raster graphics
\usepackage{graphicx}
\DeclareGraphicsExtensions{.png,.png,.jpg,.pdf}

% basic packages
\usepackage{fancyhdr}
\usepackage{color}
\usepackage{parskip}
\usepackage{float}
\usepackage{microtype}
\usepackage{url}
\usepackage{hyperref}
\usepackage{booktabs}
\usepackage{makecell}
\usepackage{multirow}
\usepackage{pbox}
\usepackage{enumitem}

% fonts
\usepackage{fontspec}
\setmainfont{DroidSans}[
  Extension   = .ttf,
  UprightFont = *,
  BoldFont    = *-Bold ]
\setmonofont{DroidSansMono}[
  Extension   = .ttf,
  UprightFont = * ]

\hypersetup{colorlinks,urlcolor=blue}
\urlstyle{same}

% page layout
\usepackage[top=1in, bottom=1in, left=0.75in, right=0.75in, includefoot, heightrounded]{geometry}

% Useful for seeing the layout as you tweak it.
%\usepackage{showframe}

\setlength{\headheight}{1.1in}
\pagestyle{fancy}
\fancyhf{}

\lhead{\includegraphics[height=1in,clip]{logo}}
\rhead{\pbox[b][1in][c]{\textwidth}{
    \small
    \textbf{Prueba realizada por:}\\
    Northwest Genomics Center\\
    University of Washington School of Medicine\\
    Genome Sciences Box 355065\\
    3720 15th Ave NE\\
    Seattle, WA 98195-5065}}

\newcommand{\PageLine}{\rule{\textwidth}{0.25mm}}

\newcommand{\link}[2]{\href{#1}{#2}\footnote{#1}}

\begin{document}

\begin{center}
\Large
\textbf{Informe de laboratorio sobre la PCR cualitativa para el SARS-CoV-2
(COVID-19) de NWGC}
\end{center}

\bigskip

\begin{description}[font=\normalfont,align=left,labelwidth=18em]
\item [Nombre del participante] \textbf{\VAR{pat_name|e}}
\item [Fecha de nacimiento del participante] \textbf{\VAR{birth_date|e}}
\item [Identificador de la muestra] \textbf{\VAR{qrcode|e}}
\item [Fecha de entrega de la muestra] \textbf{\VAR{collect_ts|e}}
\item [Fecha de entrega de los resultados]
  %- if status_code == "never-tested"
  \textbf{Nunca probado}
  %- else
  \textbf{\VAR{result_ts|e}}
  %- endif
\item [Tipo de muestra] Muestra recolectada por el participante, de la sección
  media de la nariz
\end{description}

\PageLine

Resultado de la PCR (Polymerase Chain Reaction, reacción en cadena de la
polimerasa) cualitativa para el SARS-CoV-2 (Severe acute respiratory syndrome
coronavirus 2, coronavirus 2 del síndrome respiratorio agudo grave) (COVID-19)

%- if status_code == "negative"
\textbf{No detectado}\\
%- elif status_code == "positive"
\textbf{SARS-CoV-2 detectado (positivo)}\\
%- elif status_code == "inconclusive"
\textbf{No concluyente}\\
%- elif status_code == "never-tested"
\textbf{No es posible analizar la muestra}\\
%- else
\VAR{0/0} %# poor man's raise
%- endif

\PageLine

Tenga en cuenta lo siguiente:

\begin{itemize}
\item

  \textbf{El tipo de muestra y análisis que se utilizan en este proyecto son
  parte de un programa de vigilancia de la salud pública, que se realiza en
  asociación con el condado de King y otros organismos.} El análisis que recibió
  no reemplaza la atención médica personal que brindan los proveedores de
  atención médica.

\item

  \textbf{Cadena de custodia de la muestra:} la muestra analizada fue
  recolectada y etiquetada por el participante, y, luego, enviada a la SCAN por
  medio de correo postal o un servicio de entregas. No fue recolectada,
  etiquetada y transportada por un profesional de atención médica con licencia.
  Esto puede afectar los resultados del análisis, y no se pueden excluir los
  resultados falsos negativos. Podría justificarse la realización de análisis
  clínicos de seguimiento.

\item

  \textbf{Dispositivo de recolección de la muestra:} el exudado de cornete medio
  no es un tipo de muestra aprobada por medio de una EUA (Emergency Use
  Authorization, autorización de uso de emergencia) de la Food and Drug
  Administration (FDA, Administración de Alimentos y Medicamentos) o los Centers
  for Disease Control and Prevention (CDC, Centros para el Control y la
  Prevención de Enfermedades​). Por lo tanto, es posible que los resultados que
  indican “no detectado” o “negativo” no excluyan una infección por SARS-CoV-2.

\item

  \textbf{El Washington State Department of Health (Departamento de Salud del
  Estado de Washington) puede verificar los casos con resultado ``positivo''.}
  Si tiene preguntas, debe hablar sobre este resultado con su proveedor de
  atención médica.

\item{
  \textbf{Si el resultado del análisis indica ``no detectado'', esto significa
  que no se detectó el SARS-CoV-2 en su muestra en nuestro laboratorio.}

  Ningún análisis es perfecto, por lo que, aún con un resultado negativo, existe
  una pequeña posibilidad de que usted tenga el virus SARS-CoV-2 que causa la
  COVID-19. Para este análisis, no se utiliza la forma estándar de recolección
  de muestras por medio de un proveedor de atención médica; en cambio, la
  muestra se envió por correo postal o un servicio de entregas, en lugar de
  recolectarse y transportarse en un entorno de atención médica.
  Independientemente de los resultados que obtenga en el análisis, si se siente
  enfermo, se recomienda que se quede en su casa y evite el contacto con otras
  personas. Si tiene otros problemas de salud o está muy enfermo, o si siente
  preocupación por su enfermedad, debe comunicarse con su proveedor de atención
  médica. Si los síntomas persisten, debe hablar sobre la posibilidad de
  realizarse otro análisis para detectar la COVID-19 con su proveedor de
  atención médica.
}

\item

  \textbf{Si el resultado del análisis indica ``no concluyente'',} esto
  significa que los resultados del análisis realizado en nuestro laboratorio
  para detectar el SARS-CoV-2 no son certeros.

  Ningún análisis es perfecto, por lo que existe la posibilidad de que un
  resultado no concluyente pueda significar que usted tiene el virus SARS-CoV-2
  que causa la COVID-19. Para este análisis, no se utiliza la forma estándar de
  recolección de muestras por medio de un proveedor de atención médica; en
  cambio, la muestra se envió por correo postal o un servicio de entregas, en
  lugar de recolectarse y transportarse en un entorno de atención médica.
  Independientemente de los resultados que obtenga en el análisis, si se siente
  enfermo, se recomienda que se quede en su casa y evite el contacto con otras
  personas. Si tiene otros problemas de salud o está muy enfermo, o si siente
  preocupación por su enfermedad, debe comunicarse con su proveedor de atención
  médica. Si los síntomas persisten, debe hablar sobre la posibilidad de
  realizarse otro análisis para detectar la COVID-19 con su proveedor de
  atención médica.

\item

  \textbf{Si el resultado del análisis indica ``no es posible analizar la
  muestra'',} esto significa que la cantidad de muestra no fue suficiente o que
  se produjo un error durante el proceso del análisis.

Dado que no fue posible realizar el análisis, existe la posibilidad de que usted
 tenga el virus SARS-CoV-2 que causa la COVID-19. Para este análisis, no se
 utiliza la forma estándar de recolección de muestras por medio de un proveedor
 de atención médica; en cambio, la muestra se envió por correo postal o un
 servicio de entregas, en lugar de recolectarse y transportarse en un entorno de
 atención médica. Independientemente de los resultados que obtenga en el
 análisis, si se siente enfermo, se recomienda que se quede en su casa y evite
 el contacto con otras personas. Si tiene otros problemas de salud o está muy
 enfermo, o si siente preocupación por su enfermedad, debe comunicarse con su
 proveedor de atención médica. Si los síntomas persisten, debe hablar sobre la
 posibilidad de realizarse otro análisis para detectar la COVID-19 con su
 proveedor de atención médica.


\item

  \textbf{El centro en el que se realiza el análisis informará los casos en los
  que los análisis arrojen resultados positivos (en los que se indique
  ``DETECTADO'' o ``Virus detectado'') al departamento de Salud del condado o
  estado correspondiente.} NO es necesario que informe estos resultados usted
  mismo. Si el resultado de su análisis es positivo para el nuevo coronavirus,
  cabe la posibilidad de que Salud Pública realice un seguimiento de su caso, y
  le haga preguntas sobre su enfermedad y las personas con las que pudo haber
  estado en contacto.

\item

  \textbf{Puede imprimir los resultados de su análisis o enviárselos por correo
  electrónico a su proveedor de atención médica.}

\end{itemize}

\bigskip
\textbf{Metodología de la PCR cualitativa para el SARS-CoV-2 (COVID-19):}

Este análisis está diseñado para detectar el SARS-CoV-2 por medio de la PCR con
transcripción inversa. El análisis se basa en la amplificación de los segmentos
Orf1B y S del genoma viral con control interno sobre la amplificación de la
ribonucleasa P (RNasa P) humana. El Northwest Genomics Center (NWGC) de la
University of Washington desarrolló este análisis y determinó sus
características de funcionamiento. El análisis desarrollado por este laboratorio
está validado y fue presentado ante la Food and Drug Administration de los EE. 
UU. para obtener la autorización de uso de emergencia. La evaluación
independiente de la validación del análisis por parte de la FDA se encuentra
pendiente. Este laboratorio cuenta con la certificación de las Clinical
Laboratory Improvement Amendments (CLIA, enmiendas para la mejora de
laboratorios clínicos) y está calificado para realizar análisis de laboratorio
clínicos de alta complejidad.

Para obtener más información acerca del proveedor, consulte los siguientes
enlaces proporcionados por el Washington State Department of Health:

\begin{itemize}
\item

  \link{https://www.doh.wa.gov/Portals/1/Documents/1600/coronavirus/COVIDcasepositiveSPANISH.pdf}{Pacientes
  con sospecha o confirmación de la COVID-19}

\item

  \link{https://www.doh.wa.gov/Portals/1/Documents/1600/coronavirus/COVIDexposed\_notSickYet\_ES.pdf}{Pacientes
  que estuvieron expuestos a un caso confirmado de la COVID-19}

\item

  \link{https://www.doh.wa.gov/Portals/1/Documents/1600/coronavirus/COVIDsickNotExposedSpanish.pdf}{Pacientes
  que no estuvieron expuestos y que tienen síntomas de la COVID-19}

\end{itemize}

En los sitios web del Washington State Department of Health y de los CDC se
ofrecen recursos adicionales:

\begin{itemize}
\item

  \url{https://www.doh.wa.gov/esp}

\item

  \url{https://espanol.cdc.gov/coronavirus/2019-ncov/index.html}

\end{itemize}

\end{document}
